\documentclass[11pt,a4paper]{article}
\usepackage[utf8]{inputenc}
\usepackage[margin=2.4cm]{geometry}
\usepackage{amsmath,amssymb}
\usepackage{booktabs}
\usepackage{xcolor}
\usepackage{graphicx}
\usepackage{listings}
\usepackage{enumitem}
\usepackage{microtype}
\usepackage[colorlinks=true,linkcolor=blue!50!black,urlcolor=blue!50!black]{hyperref}

\definecolor{codegray}{gray}{0.45}
\definecolor{codeblue}{rgb}{0.1,0.2,0.6}
\lstset{
  language=Python,
  basicstyle=\ttfamily\footnotesize,
  keywordstyle=\color{codeblue},
  commentstyle=\color{codegray}\itshape,
  stringstyle=\color{teal},
  showstringspaces=false,
  breaklines=true,
  frame=single,
  rulecolor=\color{black!25},
  numbers=left,
  numberstyle=\tiny\color{codegray},
  literate={θ}{$\theta$}1 {ξ}{$\xi$}1,
}

\newcommand{\session}[3]{%
  \title{\textbf{#1}\\[2pt]\large #2}
  \author{astro-lab \,--\, #3}
  \date{}
  \maketitle
  \vspace{-1.5em}
}
\newcommand{\todo}[1]{\textcolor{red!60!black}{\textbf{[#1]}}}
\usepackage{framed}
\newenvironment{keyidea}{%
  \def\FrameCommand{\color{blue!40!black}\vrule width 2.5pt\hspace{8pt}\color{black}}%
  \MakeFramed{\advance\hsize-\width\FrameRestore}}{\endMakeFramed}

\newcommand{\dd}{\mathrm{d}}
\newcommand{\pc}{\,\mathrm{pc}}
\newcommand{\Msun}{M_\odot}
\newcommand{\Rsun}{R_\odot}
\setlist[itemize]{itemsep=2pt,topsep=4pt}
\setlist[enumerate]{itemsep=3pt,topsep=4pt}
\setlength{\parindent}{0pt}
\setlength{\parskip}{6pt}

\begin{document}
\session{Session 4: Radiative Processes}{Blackbody radiation, temperature scales, spectral lines}{reading notes}

\textbf{Status:} reading notes and exercise specification. The code for this session (\texttt{04\_radiation/}) is not in the repository yet; the exercises below define what to build.

\section*{Goals}
Use the Planck function to relate a star's colour, temperature and luminosity, and fit an absorption line profile to extract a line centre, width and equivalent width.

\section{The Planck function}
The specific intensity of a blackbody at temperature $T$, per unit wavelength, is
\[
B_\lambda(T) = \frac{2hc^2}{\lambda^5}\,\frac{1}{\exp\!\left(\dfrac{hc}{\lambda k_BT}\right)-1}.
\]
Two limits are worth deriving yourself:
\begin{itemize}
\item \textbf{Rayleigh--Jeans} ($hc/\lambda k_BT\ll1$): $B_\lambda\approx 2ck_BT/\lambda^4$.
\item \textbf{Wien} ($hc/\lambda k_BT\gg1$): $B_\lambda\approx (2hc^2/\lambda^5)\,e^{-hc/\lambda k_BT}$.
\end{itemize}

\section{Wien and Stefan--Boltzmann}
Setting $\dd B_\lambda/\dd\lambda=0$ gives Wien's displacement law,
\[
\lambda_{\max}T = b \approx 2.898\times10^{-3}\ \mathrm{m\,K}.
\]
Integrating over all wavelengths and solid angle gives the flux $F=\sigma T^4$ with $\sigma\approx5.670\times10^{-8}\ \mathrm{W\,m^{-2}\,K^{-4}}$. For a spherical star,
\[
L = 4\pi R^2\sigma T_{\mathrm{eff}}^4 .
\]
Using $L_\odot=3.828\times10^{26}$ W and $R_\odot=6.957\times10^{8}$ m gives $T_{\mathrm{eff},\odot}\approx 5772$ K, a good check for your code.

\begin{keyidea}
\textbf{Why $L\propto R^2T^4$ matters.} A red giant and a white dwarf can share a colour but differ in luminosity by many orders of magnitude, purely through $R$. This is the geometric reason the HR diagram has distinct branches.
\end{keyidea}

\section{Spectral lines}
Line profiles come from three broadening mechanisms:
\begin{itemize}
\item \textbf{Natural}: Lorentzian, from the finite lifetime of the level.
\item \textbf{Doppler} (thermal and turbulent): Gaussian, $\sigma_D=\dfrac{\lambda_0}{c}\sqrt{\dfrac{k_BT}{m}}$ for thermal motion.
\item \textbf{Pressure}: Lorentzian, from collisions.
\end{itemize}
Gaussian and Lorentzian combine into a \textbf{Voigt} profile (their convolution), usually evaluated with the Faddeeva function (\texttt{scipy.special.wofz}). The \textbf{equivalent width}
\[
W_\lambda = \int\left(1-\frac{F_\lambda}{F_c}\right)\dd\lambda
\]
is the width of a fully absorbed rectangle with the same area; it measures line strength independent of resolution.

\section{Implementation plan}
\begin{enumerate}
\item Write \texttt{planck(lam, T)} with units; test the Wien peak numerically against $b/T$ for $T=3000$, $5772$, $10000$ K.
\item Integrate the Planck function numerically and check $\sigma T^4/\pi$ for the intensity integral.
\item Generate a synthetic spectrum: Planck continuum $\times$ a Voigt absorption line plus Gaussian noise.
\item Fit with \texttt{scipy.optimize.curve\_fit}; report best-fit centre, width and uncertainties from the covariance matrix.
\item Compare the fitted line depth with the equivalent width computed by direct integration.
\end{enumerate}

\section{Exercises}
\begin{enumerate}
\item Compute the peak wavelength for the Sun, Betelgeuse ($\sim3600$ K) and Rigel ($\sim12\,000$ K) and say which lie in the visible range.
\item Derive the Rayleigh--Jeans limit from the Planck function.
\item Estimate the thermal Doppler width of H$\alpha$ (656.28 nm) at 6000 K and compare with typical observed widths.
\item Show how the fitted width uncertainty scales with the noise level and with the number of pixels across the line.
\end{enumerate}

\section*{References}
\begin{itemize}
\item Rybicki \& Lightman, \textit{Radiative Processes in Astrophysics}, ch.\ 1 and 10
\item Gray, \textit{The Observation and Analysis of Stellar Photospheres}
\item Carroll \& Ostlie, \textit{An Introduction to Modern Astrophysics}, ch.\ 3 and 8
\end{itemize}
\end{document}
